\documentclass[a4paper,10pt]{article}
\usepackage[left=1in, right=1in, top=1in, bottom=1in]{geometry}
\usepackage{enumitem}
\usepackage{titlesec}
\usepackage{parskip}
\usepackage{xcolor}
\usepackage{fontawesome5}
\usepackage[hidelinks]{hyperref}
\usepackage{fancyhdr}
\usepackage{needspace}

% Define a cor verde oliva
\definecolor{verdeolivo}{RGB}{85,107,47}

\titleformat{\section}{\large\bfseries\color{verdeolivo}}{}{0em}{}[\titlerule]
\titleformat{\subsection}[runin]{\bfseries\color{verdeolivo}}{}{0em}{}[:]

\pagestyle{fancy}
\fancyhf{}
\fancyfoot[L]{Última atualização: \today}

\begin{document}

\begin{center}
    \textbf{\Huge Antonio Neto} \\
    \vspace{2mm}
    \faIcon{briefcase} AI Engineer $|$ Estudante de Ciência da Computação
\end{center}

\vspace{4mm}
\begin{center}
    \faIcon{graduation-cap} \href{https://lattes.cnpq.br/6284651096866272}{Lattes} \hspace{1.5cm}
    \faIcon{github} \href{https://github.com/nietus}{Github} \hspace{1.5cm}
    \faIcon{envelope} \href{mailto:antonio@deeptera.com}{Email} \hspace{1.5cm}
    \faIcon{linkedin} \href{https://www.linkedin.com/in/antonioniet/}{Linkedin} \hspace{1.5cm}
    \faIcon{globe} \href{https://antonio.niet.com.br}{Website}
\end{center}

\space
\section*{\faIcon{user} Perfil Profissional}
Sou um AI Engineer brasileiro e estudante de Ciência da Computação na PUC Minas, trabalhando com IA aplicada, otimização, infraestrutura e problemas de logística. Minha atuação em pesquisa e desenvolvimento abrange sistemas multiagentes baseados em LLM, metaheurísticas e aprendizado de máquina aplicados a cenários industriais e logísticos, com cinco artigos aceitos. Também estudo Mandarim no Instituto Confúcio da UFMG. Sou altamente focado, detalhista e comprometido com tudo que faço.

\section*{\faIcon{building} Experiência Profissional}

\subsection*{Deeptera \hfill abr 2026 -- Presente}
\textit{Artificial Intelligence Engineer} \\
Atuação remota em meio período com pesquisa e desenvolvimento
\begin{itemize}
    \item Pesquisa e desenvolvimento em problemas de otimização e sistemas multiagentes.
    \item Concepção de arquiteturas para sistemas inteligentes integrados a softwares convencionais.
    \item Autoria de artigos aceitos no SBPO e no SIMPEP sobre otimização e sistemas multiagentes.
\end{itemize}

\subsection*{Deeptera \hfill ago 2024 -- abr 2026}
\textit{Estágio em computação e dados} \\
Pesquisa e desenvolvimento relacionados a Inteligência Artificial
\begin{itemize}
    \item Artigo sobre sistemas multiagentes aprovado no SBAI.
    \item Participação em projetos financiados pela FAPEMIG.
    \item Arquitetura de bancos de dados e backend em Django com infraestrutura Azure e Azure DevOps.
    \item Prestação de serviços de otimização e backend para a VLI.
    \item Desenvolvimento e pesquisa em sistemas multiagentes, com aplicação de algoritmos de RAG, ferramentas personalizadas e serviços escaláveis.
    \item Aplicação de algoritmos genéticos na resolução de problemas de otimização.
\end{itemize}

\subsection*{Hackatruck \hfill 2023}
\textit{Participante do Projeto} \\
Desenvolvi aplicações iOS como parte de um projeto em grupo de um mês.
\begin{itemize}
    \item Colaborei com membros da equipe para projetar e implementar funcionalidades do aplicativo.
    \item Utilizei Swift, Node Red e Azure para desenvolvimento, lidando com várias APIs.
\end{itemize}

\section*{\faIcon{book} Publicações}
\begin{itemize}[label=\faIcon{file-alt}]
    \needspace{4\baselineskip}\item \textbf{Uma Abordagem Baseada em Atenção para o Problema de Coletas e Entregas Dinâmicas.} A. Neto, R. D. Macedo, A. H. Duque, S. J. de Almeida. LVIII SBPO, 2026. \\
    Modelo de Atenção treinado com Aprendizado por Reforço e combinado com VNS para o Problema de Coletas e Entregas Dinâmicas, avaliado no benchmark ICAPS 2021.
    \needspace{4\baselineskip}\item \textbf{DeepFactory: Sistema Multiagentes para Gestão e Otimização de Produção Industrial.} R. D. Macedo, A. Neto, A. H. Duque, S. J. de Almeida. LVIII SBPO, 2026. \\
    Plataforma multiagente orientada a estágios, integrando sequenciamento via MIP, previsão de demanda, manutenção preditiva e controle de qualidade por visão computacional.
    \needspace{4\baselineskip}\item \textbf{Otimização do Sequenciamento de Operações de Carregamento e Descarregamento de Navios Graneleiros com Restrição de Trim via Algoritmo Genético.} A. Neto, R. D. Macedo, C. S. Almeida et al. (com a VLI Multimodal). LVIII SBPO, 2026. \\
    Algoritmo genético com restrição de trim validado em 112 planos operacionais reais de uma empresa de logística multimodal.
    \needspace{4\baselineskip}\item \textbf{Roteamento Multi-Período de Drones com Janelas de Tempo e Disponibilização Gradual de Estações de Recarga.} A. S. C. Neto, Z. K. G. do Patrocínio. SIMPEP, 2026. \\
    Trabalho de Conclusão de Curso: MP-EAVNS, uma Busca em Vizinhança Variável sensível à energia para uma extensão multi-período do E-VRPTW.
    \needspace{4\baselineskip}\item \textbf{Proposta de um Sistema Multiagentes baseado em LLM: Integração em um Software de Gerenciamento e Otimização de Carregamento e Descarregamento de Navios.} A. Neto, C. S. Almeida, R. D. Macedo, A. H. Duque, S. J. de Almeida. SBAI/SBSE, 2025. \\
    Sistema multiagente hierárquico com LLM (Manager, Helper, Optimizer e Data Analytics) integrado a um software convencional de carregamento e descarregamento de navios.
\end{itemize}

\section*{\faIcon{graduation-cap} Educação}

\subsection*{PUC Minas \hfill 2023 -- Presente}
\textit{Bacharelado em Ciência da Computação} \\
Tenho trabalhado com várias tecnologias na PUC e me especializado em aprendizado de máquina, ciência de dados e tecnologias de IA fora da faculdade. Meu Trabalho de Conclusão de Curso trata do roteamento multi-período de drones com janelas de tempo e disponibilização gradual de estações de recarga, orientado pelo Prof. Zenilton K. G. do Patrocínio.

\subsection*{Instituto Confúcio (UFMG) \hfill 2024 -- Presente}
\textit{Estudante de Mandarim (HSK 4)} \\
Estudo Mandarim por gostar da cultura e pela sua relevância na área de IA.

\section*{\faIcon{cogs} Principais Habilidades Técnicas}
\begin{center}
    \begin{itemize}[label=\faIcon{check}, itemsep=-3pt]
        \item \textbf{Linguagens de Programação:} Python, C, TypeScript
        \item \textbf{Sistemas de Agentes:} LangChain, LangGraph, RAG, MCP
        \item \textbf{Aprendizado de Máquina:} PyTorch, Scikit-learn, XGBoost
        \item \textbf{Ciência de Dados:} Pandas, NumPy, Matplotlib, Seaborn
        \item \textbf{Backend:} Django, FastAPI, APIs REST
        \item \textbf{Bancos de Dados:} PostgreSQL, MySQL
        \item \textbf{DevOps:} Azure DevOps, Docker
        \item \textbf{Controle de Versão:} Git, GitHub
        \item \textbf{Otimização:} AMPL, pyomo, solvers GLPK/SCIP, metaheurísticas (AG, VNS, ALNS, GRASP)
        \item \textbf{Cloud:} Azure, Google Cloud
    \end{itemize}
\end{center}

\section*{\faIcon{certificate} Certificações}
\begin{itemize}[label=\faIcon{certificate}]
    \item Práticas de Cloud Services usando Swift com ênfase em Serviços Cognitivos
    \item Algoritmos e POO com Swift, JavaScript e RESTful APIs
    \item Agricultura 4.0: Otimizando com Inteligência Artificial
\end{itemize}

\section*{\faIcon{language} Idiomas}
\begin{itemize}[label=\faIcon{globe}]
    \item Português - Nativo
    \item Inglês - Fluente
    \item Mandarim - HSK 4
\end{itemize}

\end{document}
